\section{Biological background}
%==========================================================================

\subsection{Cell and gene therapy in rare disease}

Most rare diseases are monogenic: a single gene is defective and, as a consequence, a
single protein is absent or non-functional. This makes them structurally different from
common diseases, where hundreds of loci contribute small effects. It also makes them
tractable: if one knows which protein is missing, the therapeutic objective is simply to
restore it.

Cell and gene therapies do exactly this. Rather than administering a drug that modulates
a pathway, they supply the genetic instruction for the missing protein, either by editing
the patient's genome, by delivering a functional copy of the gene, or by delivering
messenger RNA that the patient's own cells translate. The last of these is transient: the
mRNA is degraded within days, so the therapy must be re-administered indefinitely.

The transience is a limitation clinically and an advantage mathematically. A permanent
genomic correction is a single irreversible intervention. A repeated mRNA dose is a
\emph{manipulable input}, applied on a known schedule, with a measurable response between
administrations. It is this structure that makes the present problem a control problem.

\subsection{Propionic acidemia}

Propionic acidemia (PA) is an autosomal recessive disorder of propionate metabolism.
Four amino acids (isoleucine, valine, methionine, threonine), odd-chain fatty acids, and
propionate produced by intestinal flora all converge on a common intermediate,
propionyl-CoA. Under normal conditions this intermediate is carboxylated to
methylmalonyl-CoA and enters the tricarboxylic acid cycle as succinyl-CoA.

The enzyme performing this step is propionyl-CoA carboxylase (PCC), a biotin-dependent
mitochondrial dodecamer $\alpha_6\beta_6$~\cite{pcc_structure}. % C01
The $\alpha$ subunit is encoded by \emph{PCCA}
and carries the biotin cofactor; the $\beta$ subunit is encoded by \emph{PCCB} and
performs the carboxyl transfer. Both subunits are required, which is why pathogenic
variants in either gene produce the same disease.

\begin{adjustwidth}{-1.3cm}{-1.3cm}
\begin{figure}[htbp]
\centering

%,,,,,,,,,,,,,,,,,,,,,,,,,,,,,,,,,,,,,,,,,,,,,,,,, (a) the pathway
\begin{minipage}[t]{0.57\linewidth}
\centering
\resizebox{\linewidth}{!}{%
\begin{tikzpicture}[
  box/.style={draw, rounded corners=2pt, align=center, inner sep=5pt,
              font=\small, minimum height=9mm},
  src/.style={box, fill=soft},
  mid/.style={box, fill=white, font=\small\bfseries},
  bad/.style={box, fill=red!8, draw=red!50},
  esc/.style={box, fill=blue!6, draw=blue!40},
  ok/.style={box, fill=green!8, draw=green!45!black!50},
  ar/.style={-{Stealth[length=2.2mm]}, thick},
]

%, sources, left column
\node[src] (aa)  at (0,1.6)  {Ile, Val, Met, Thr\\{\scriptsize dietary and catabolic}};
\node[src] (fa)  at (0,0)    {odd-chain fatty acids};
\node[src] (gut) at (0,-1.4) {gut flora};

%, central pool
\node[mid] (pcoa) at (4.9,0) {propionyl-CoA};

\draw[ar] (aa)  -- (pcoa);
\draw[ar] (fa)  -- (pcoa);
\draw[ar] (gut) -- (pcoa);

%, blocked enzymatic route, right
\node[ok] (mmcoa) at (10.4,0) {methylmalonyl-CoA\\$\rightarrow$ TCA cycle};
\draw[ar] (pcoa) -- (mmcoa);
\node[font=\scriptsize, align=center] at (7.65,0.75) {PCC \; $\alpha_6\beta_6$};
\node[font=\scriptsize, align=center, red!70!black] at (7.65,-0.62) {capacity $\theta$};
\draw[very thick, red!70!black] (7.65,-0.34) -- (7.65,0.34);

%, escape route up
\node[esc] (carn) at (10.1,3.0) {propionylcarnitine\\{\scriptsize urinary excretion}};
\draw[ar, blue!55!black] (pcoa) to[out=60,in=190] (carn);
\node[font=\scriptsize, align=center, blue!55!black, fill=white, inner sep=1pt]
      at (6.7,1.95) {$+$ carnitine\\{\scriptsize consumed, not catalytic}};

%, escape route down
\node[bad] (tox) at (10.1,-3.0) {2-methylcitrate\\3-hydroxypropionate\\{\scriptsize measured biomarkers}};
\draw[ar, red!60!black] (pcoa) to[out=-55,in=180] (tox);

%, sequestration
\node[bad, text width=34mm] (coa) at (4.1,-3.0)
     {CoA sequestration\\{\scriptsize $\Rightarrow$ PDH, $\beta$-oxidation, TCA impaired}};
\draw[ar, red!60!black] (pcoa) -- (coa);

\end{tikzpicture}}

\vspace{0.5em}
{\small\textbf{(a)} the pathway}
\end{minipage}\hfill
%,,,,,,,,,,,,,,,,,,,,,,,,,,,,,,,,,,,,,,,,,,,,,,,,, (b) the model
\begin{minipage}[t]{0.40\linewidth}
\centering
\resizebox{\linewidth}{!}{%
\begin{tikzpicture}[
  box/.style={draw, rounded corners=2pt, align=center, inner sep=5pt,
              font=\small, minimum height=9mm},
  pool/.style={box, fill=soft, font=\small\bfseries, minimum width=24mm},
  enzy/.style={box, fill=green!8, draw=green!45!black!50, minimum width=26mm},
  altc/.style={box, fill=blue!6, draw=blue!40, minimum width=26mm},
  ar/.style={-{Stealth[length=2.2mm]}, thick},
]

\node[pool] (m) at (0,0) {$m$\\[1pt]\scriptsize\mdseries accumulated\\metabolite};
\node[font=\small] (in) at (-2.7,0) {$s(t)$};
\draw[ar] (in) -- (m);
\node[font=\scriptsize, align=center] at (-1.95,0.42) {load};

\node[enzy] (enz) at (-1.85,-2.5) {$\theta\, g(m)$\\{\scriptsize saturating}};
\node[altc] (alt) at ( 1.85,-2.5) {$k_e\, m$\\{\scriptsize linear}};
\draw[ar] (m.south) -- ++(0,-0.45) -| (enz.north);
\draw[ar] (m.south) -- ++(0,-0.45) -| (alt.north);

\node[font=\scriptsize] (o1) at (-1.85,-4.1) {eliminated};
\node[font=\scriptsize] (o2) at ( 1.85,-4.1) {eliminated};
\draw[ar] (enz) -- (o1);
\draw[ar] (alt) -- (o2);

\node[font=\small, align=center, blue!55!black] (carn) at (4.3,-2.5)
     {carnitine\\$u_c(t)$};
\draw[ar, blue!55!black] (carn) -- (alt);

\end{tikzpicture}}

\vspace{0.5em}
{\small\textbf{(b)} the model}
\end{minipage}

%,,,,,,,,,,,,,,,,,,,,,,,,,,,,,,,,,,,,,,,,,,,,,,,,, correspondence
\vspace{1.1em}
{\footnotesize
\begin{tabular}{@{}p{0.44\linewidth} p{0.16\linewidth} p{0.34\linewidth}@{}}
\toprule
\textbf{(a) biology} & \textbf{(b) model} & \textbf{status of the identification} \\
\midrule
amino acids, gut flora, protein catabolism & $s(t)$ &
input; partially recorded \\
intracellular propionyl-CoA & \emph{not modelled} &
the enzyme's true substrate, never measured \\
plasma 2-MC, 3-HP, C3 acylcarnitine & $m$ \ \textbf{[proxy]} &
the only observable; a downstream marker, not the substrate \\
PCC, its $V_{\max}$ & $\theta$ &
the unknown the decision needs \\
carnitine conjugation and renal excretion & $k_e$ &
assumed population-level \\
CoA sequestration & \emph{absent in A and B} &
this is exactly what Model C reintroduces \\
\bottomrule
\end{tabular}}

\caption{(a) Propionate handling. Substrate converges on propionyl-CoA from diet,
endogenous protein catabolism, fatty acids and gut flora. The enzymatic route through PCC
is impaired; residual capacity is denoted $\theta$. Two escape routes remain: conjugation
with carnitine, which consumes a finite cofactor, and conversion to the toxic metabolites
that are measured clinically. Accumulated propionyl-CoA sequesters free coenzyme A, which
is a shared and limited resource. (b) The same system as the block diagram that the
rest of this note analyses, containing only quantities the model has. A single pool $m$ is
filled by the load $s(t)$ and drained by two channels, one saturating and one linear, with
carnitine entering as an input on the second. The table gives the correspondence, and
names the two places where the reduction discards something.}
\label{fig:pathway}
\end{figure}
\end{adjustwidth}

\paragraph{What the reduction discards.} Panel (b) is not panel (a) with the labels
removed. Three things are lost in passing from one to the other, and they are stated here
rather than left for a reader to find.

\begin{enumerate}[leftmargin=*, itemsep=3pt]
\item \textbf{The state $m$ is not the enzyme's substrate.} The substrate of PCC is
intracellular, mitochondrial propionyl-CoA, and it is never measured. What is measured is
plasma 2-methylcitrate, 3-hydroxypropionate and C3 acylcarnitine, which are downstream
products in a different compartment. Writing the Michaelis-Menten term as $\theta g(m)$
with $m$ the plasma marker asserts a map between the two, and that map is an assumption,
not a measurement. It is the single largest liberty taken in this note. Everything
downstream inherits it, including the numerical value of $K_m$, which is therefore not the
in vitro Michaelis constant of PCC.
\item \textbf{Urea cycle inhibition and hyperammonemia are not modelled.} Propionyl-CoA
inhibits N-acetylglutamate synthase and impairs ammonia elimination through a route that
does not pass through $m$ at all. Hyperammonemia is a major part of what makes a
decompensation dangerous and it is absent from the model. A patient can therefore be in
serious trouble for a reason the state variable cannot express.
\item \textbf{CoA sequestration is absent from Models A and B.} The feedback in the
lower-left of panel (a), whereby accumulated propionyl-CoA consumes a shared cofactor and
degrades the cell's general metabolic capacity, has no counterpart in a model whose
clearance is a fixed function of $m$. This is not an oversight to be apologised for: it is
precisely the mechanism Model C reintroduces, and the reason Model C is the only one of the
three in which a crisis fails to resolve by itself.
\end{enumerate}

These are limits of the reduction, not defects of the biology, and each is a statement
that can be checked. The first is checkable against any study relating plasma markers to
intracellular propionyl-CoA. The second bounds what the model may be used to predict. The
third is answered by Model C.

\subsection{Why the block is damaging}

Two mechanisms, both visible in Figure~\ref{fig:pathway}.

\paragraph{Coenzyme A sequestration.} Coenzyme A is a finite, shared carrier. Accumulated
propionyl-CoA holds a growing fraction of the cellular pool, leaving less free CoA for
pyruvate dehydrogenase, $\beta$-oxidation and the TCA cycle. Energy production therefore
fails not because any of those enzymes is defective, but because a shared resource has
been consumed. This is the origin of the secondary mitochondrial dysfunction
characteristic of the disease.

\paragraph{Urea cycle inhibition.} Propionyl-CoA inhibits N-acetylglutamate synthase, the
activator of the urea cycle. Ammonia elimination is impaired and hyperammonemia develops
independently of the acidosis, through a separate route from the same cause.

\subsection{Clinical course}

Patients live in a fragile steady state maintained by protein restriction, carnitine
supplementation and, frequently, intestinal antibiotics. Toxic metabolites are permanently
elevated but tolerated.

The dangerous events are \emph{metabolic decompensations}. They are triggered by anything
that raises substrate load: infection, fasting, surgery, trauma, rapid growth, or a
dietary error. The dominant mechanism is counter-intuitive. During intercurrent illness the
patient stops eating, the body enters catabolism, and skeletal muscle protein is degraded
for energy. Since muscle contains the four precursor amino acids, \textbf{the propionate
load rises sharply precisely when the patient is eating nothing}. The source is internal
and cannot be controlled by diet.

The resulting crisis combines metabolic acidosis and hyperammonemia, progressing to
lethargy, coma and, without intervention, death. Treatment is directed at each term of the
balance: intravenous glucose to suppress catabolism, total protein withdrawal, intravenous
carnitine, ammonia scavengers, and haemofiltration in refractory cases. Most crises
resolve, but each leaves irreversible neurological damage, with the basal ganglia
particularly vulnerable.

Reported event rates are approximately $0.5$ decompensations per patient-year averaged
across ages~\cite{mde_claims_us}, % C05
but strongly age-dependent: around $1.4$ per patient-year between ages 3 and 6, falling to
$0.1$ between 13 and 17~\cite{mde_claims_us}. % C06
The profile is \emph{U-shaped} rather than monotonically decreasing. The same series reports
the rate rising again in adulthood, to $0.2$ per patient-year~\cite{mde_claims_us}, % C07
and a separate study of United States electronic records reports a substantially higher
adult rate of $0.72$ per patient-year. % C31
The two disagree, and no adult rate should be quoted without saying which study produced it.
Critically, roughly $70\%$ of patients experience no event in a given observation window, so
the event distribution is strongly bimodal~\cite{natural_history}. % C08

\subsection{The Saudi context}

Consanguinity rates in Saudi Arabia are among the highest worldwide, producing founder
variants at frequencies not seen in outbred populations. Propionic acidemia ranks among the
four most frequent disorders in Saudi carrier screening, and was among the most prevalent
inherited metabolic disorders in a Riyadh newborn-screening cohort of approximately 41{,}000
($\approx 1{:}10{,}000$)~\cite{nbs_riyadh_kfmc}. % C10

A published series of 84 Saudi patients identified 16 distinct variants (7 in \emph{PCCA},
9 in \emph{PCCB})~\cite{saudi84}. % C11
Two features of that series matter here:

\begin{itemize}[leftmargin=*, itemsep=2pt]
\item The variant \emph{PCCA} c.425G$>$A (p.Gly142Asp) accounts for $70.2\%$ of families and
      is associated with a severe, neonatal-onset phenotype~\cite{saudi84}. % C12
\item \textbf{All families but one were homozygous}~\cite{saudi84}, the exception being a
      single compound heterozygous family. % C13
      Compound heterozygosity, which
      elsewhere confounds genotype--phenotype correlation because two different variants
      contribute jointly (and, for a multimeric enzyme, non-additively through intragenic
      complementation), is essentially absent.
\end{itemize}

The second point is the methodological reason to work in this population: the map from
genotype to residual enzymatic capacity is unconfounded~\cite{saudi84}. % C13
That unconfounding is bought at a price in generalisability, which should be stated in the
same breath. The series is $81\%$ \emph{PCCA} (68 of 84 patients) against $19\%$
\emph{PCCB}, where the two genes contribute roughly equally to propionic acidemia
worldwide~\cite{saudi84,genereviews_pa}. % C29
Anything estimated here is estimated predominantly on \emph{PCCA} biology. Obligate carrier parents, who
retain approximately half of normal activity and are asymptomatic, provide a natural upper
anchor.

\subsection{Why the therapy makes $\theta$ a decision variable}

An investigational LNP-formulated mRNA therapy encoding both PCC subunits is in clinical
development, administered intravenously every two weeks~\cite{trial_phase12}. % C15
Reported phase 1/2 results indicate a $70\%$ reduction in decompensation risk, but the
qualifiers belong with the number: the study is single-arm, and the reduction is computed
within the eight participants who reported an event in the twelve months before
treatment~\cite{trial_phase12}. % C16
Dose selection was performed by population PK/PD modelling using \emph{body weight} as the
covariate and \emph{biomarker reduction} as the endpoint, yielding a single recommended
regimen applied uniformly~\cite{dose_selection_poster,pkpd_translational}. % C17

Genotype is not a covariate in that analysis~\cite{dose_selection_poster}. % C18
Since the Saudi literature documents
genotype-dependent severity, the question of whether the uniform regimen is adequate for
every patient is open.

This is the reason the residual capacity $\theta$ is treated here as a decision variable
rather than a descriptive one. Before the therapy, knowing a patient's $\theta$ changes
nothing that can be done for them: diet, carnitine and antibiotics are titrated against the
biomarker, and the capacity behind it is of academic interest. With the therapy, $\theta$
is the quantity that says how much enzyme has to be supplied and how often, so the same
number becomes the input to a dosing decision. The rest of this note asks how much of it
can be known before that decision has to be made.

\paragraph{The cost of an unnecessary dose is not only financial.} It is tempting to read
the question as one of expenditure, with the worst case of a uniform regimen being that
some patients are dosed more often than they need. The phase 1/2 safety data do not support
that reading. Infusion-related reactions were reported in $11$ of $43$ doses
($25.6\%$). One participant who was positive for anti-PEG antibodies at baseline showed a
fourfold rise in titre before the third dose, which is the pattern that precedes accelerated
clearance of PEGylated particles and loss of efficacy. Every dose that a patient did not
need is an additional exposure of an immunogenic construct, and the anticipated consequence
of accumulated exposure is that the therapy stops working in the patient who needs it.
Knowing $\theta$ before dosing therefore protects the therapy itself, not merely the
budget, and it is one of the arguments for individualising a schedule rather than
shortening it uniformly.

\vspace{0.5em}
\hrule
\vspace{1em}

%==========================================================================
